% UnsolvedMath ID 3800003; source catalogue code AMR-037-0003
% !TeX program = lualatex
% Definition and sources only; this file is not an LLM solution attempt.
\documentclass[11pt,a4paper]{article}
\usepackage[margin=25mm]{geometry}
\usepackage{fontspec}
\setmainfont{lmroman10-regular.otf}[BoldFont=lmroman10-bold.otf,
 ItalicFont=lmroman10-italic.otf,BoldItalicFont=lmroman10-bolditalic.otf]
\usepackage{amsmath,amssymb,mathtools}
\usepackage{unicode-math}
\setmathfont{latinmodern-math.otf}
\AtBeginDocument{\let\setminus\smallsetminus}
\usepackage{xurl,hyperref}
\hypersetup{colorlinks=true,urlcolor=blue}
\setcounter{secnumdepth}{0}
\setlength{\parindent}{0pt}
\setlength{\parskip}{5pt}
\setlength{\emergencystretch}{3em}
\begin{document}
\section*{Degenerate facets of polytopes}\label{3800003}
{\small UnsolvedMath ID 3800003 · AMR-037-0003 · Geometry · AMR Open Problem Lists}
\subsection{Definitions and mathematical statement}
A convex $d$-polytope is the convex hull of finitely many points in
$\mathbb R^d$ with nonempty interior. Its vertices are its zero-dimensional
faces, and a facet is a face of dimension $d-1$. A facet is degenerate in
this problem when it has more than $d$ vertices; equivalently, it is not a
simplex. This use of ``degenerate'' does not mean that the whole polytope
fails to be full-dimensional.

For $n\ge d+1$ define
\[
 D_d(n)=\max_{P:\ f_0(P)=n}
 \#\{F:F\text{ is a facet of }P,\ f_0(F)>d\}.
\]
Determine $D_d(n)$, or sharp asymptotic bounds for fixed $d$ as $n\to\infty$.
The interesting gap in the historical formulation concerns even dimensions.
The source also singles out quasisimplicial polytopes, whose
$(d-2)$-dimensional faces are simplices.

The subsidiary four-dimensional question asking whether $D_4(n)\ge2n$ for
some $n$ has already been answered positively by the neighborly cubical
constructions reported in the source's update. It is therefore not posed
here as a fresh open conjecture. The definition records the remaining
extremal-count question, including its quasisimplicial restriction.
\subsection{Short English statement}
How many nonsimplicial facets can a convex polytope with a prescribed
number of vertices have? The old threshold of twice as many facets as
vertices in dimension four is already surpassed.
\subsection{Sources}
\begin{itemize}
\item[S1] ULAM.AI and the UnsolvedMath Contributors, supplied problems.json, record 3800003 (AMR-037-0003), AMR Open Problem Lists. Catalogue identity and source transcription; CC BY 4.0.
\url{https://huggingface.co/datasets/ulamai/UnsolvedMath}.
\item[S2] Erickson - Open Problems in Computational Geometry, Degenerate facets of polytopes. Original source indexed by AMR-037-0003.
\url{https://jeffe.cs.illinois.edu/open/}.
\item[S3] Jeff Erickson, Open Combinatorial Problems, degenerate facets and the 1999 update.
\url{https://jeffe.cs.illinois.edu/open/comb.html}.
\end{itemize}
\end{document}
